% Môn: Vật lý
% Lớp: 12
% Chương: Chương I. Vật lí nhiệt
% Bài: L12C1 Bài 6. Nhiệt hóa hơi riêng · Bài toán công suất, hiệu suất và thời gian nóng chảy
% ===== Câu 1 =====
% ID: L12C1B6-01-DS
% Mức: NB
\dangbt{Bài toán công suất, hiệu suất và thời gian nóng chảy}
\begin{ex}
% ID: L12C1B6-01-DS
% Mức: NB
Xét các phát biểu sau về dạng “Bài toán công suất, hiệu suất và thời gian nóng chảy” (bộ số 4).
\choiceTF
{\True H.P.t=λm khi mẫu ở nhiệt độ nóng chảy và bỏ qua phần khác.}
{Thời gian nóng chảy không phụ thuộc khối lượng.}
{\True Khi giải dạng “Bài toán công suất, hiệu suất và thời gian nóng chảy”, cần xác định đúng trạng thái đầu, trạng thái cuối và quy ước dấu hoặc điều kiện áp dụng.}
{Có thể bỏ qua mọi dữ kiện về khối lượng, nhiệt độ và hiệu suất trong mọi bài toán thuộc dạng này.}
\loigiai{
\begin{itemchoice}
\item (Đúng) H.P.t=λm khi mẫu ở nhiệt độ nóng chảy và bỏ qua phần khác. (Vì): H.P.t=λm khi mẫu ở nhiệt độ nóng chảy và bỏ qua phần khác. b) S: Thời gian nóng chảy không phụ thuộc khối lượng. c) Đ: Khi giải dạng “Bài toán công suất, hiệu suất và thời gian nóng chảy”, cần xác định đúng trạng thái đầu, trạng thái cuối và quy ước dấu hoặc điều kiện áp dụng. d) S: Có thể bỏ qua mọi dữ kiện về khối lượng, nhiệt độ và hiệu suất trong mọi bài toán thuộc dạng này. Kết luận dựa trên định nghĩa, điều kiện áp dụng và quy ước trong SGK KNTT
\item (Sai) Thời gian nóng chảy không phụ thuộc khối lượng.
\item (Đúng) Khi giải dạng “Bài toán công suất, hiệu suất và thời gian nóng chảy”, cần xác định đúng trạng thái đầu, trạng thái cuối và quy ước dấu hoặc điều kiện áp dụng.
\item (Sai) Có thể bỏ qua mọi dữ kiện về khối lượng, nhiệt độ và hiệu suất trong mọi bài toán thuộc dạng này.
\end{itemchoice}
}
\end{ex}

% ===== Câu 2 =====
% ID: L12C1B6-02-TL
% Mức: NB
\begin{bt}
% ID: L12C1B6-02-TL
% Mức: NB
Trình bày cách nhận biết và phương pháp giải dạng “Bài toán công suất, hiệu suất và thời gian nóng chảy”. Phân tích nhận định sau và sửa lại nếu sai: “Thời gian nóng chảy không phụ thuộc khối lượng.”
\loigiai{
Trước hết xác định hiện tượng, trạng thái và các đại lượng đã cho. Nêu định nghĩa hoặc hệ thức phù hợp, đổi về đơn vị SI, áp dụng đúng điều kiện và kiểm tra ý nghĩa kết quả. Nhận định cần đối chiếu: H.P.t=λm khi mẫu ở nhiệt độ nóng chảy và bỏ qua phần khác. Vì vậy nhận định đã cho sai; phát biểu đúng là: H.P.t=λm khi mẫu ở nhiệt độ nóng chảy và bỏ qua phần khác.
}
\end{bt}

% ===== Câu 4 =====
% ID: L12C1B6-03-TN
% Mức: NB
\begin{ex}
% ID: L12C1B6-03-TN
% Mức: NB
Câu 1 thuộc dạng “Bài toán công suất, hiệu suất và thời gian nóng chảy”. Phát biểu nào sau đây đúng?
\choice
{Thời gian nóng chảy không phụ thuộc khối lượng.}
{Nhiệt lượng và nhiệt độ là hai đại lượng hoàn toàn giống nhau.}
{Mọi quá trình nhiệt học đều không phụ thuộc điều kiện thực hiện.}
{\True H.P.t=λm khi mẫu ở nhiệt độ nóng chảy và bỏ qua phần khác.}
\loigiai{
Phát biểu đúng là: H.P.t=λm khi mẫu ở nhiệt độ nóng chảy và bỏ qua phần khác. Các phương án còn lại trái với mô hình, định nghĩa hoặc hệ thức nhiệt học tương ứng.
}
\end{ex}

% ===== Câu 5 =====
% ID: L12C1B6-04-DS
% Mức: TH
\begin{ex}
% ID: L12C1B6-04-DS
% Mức: TH
Xét các phát biểu sau về dạng “Bài toán công suất, hiệu suất và thời gian nóng chảy” (bộ số 1).
\choiceTF
{Thời gian nóng chảy không phụ thuộc khối lượng.}
{\True Khi giải dạng “Bài toán công suất, hiệu suất và thời gian nóng chảy”, cần xác định đúng trạng thái đầu, trạng thái cuối và quy ước dấu hoặc điều kiện áp dụng.}
{Có thể bỏ qua mọi dữ kiện về khối lượng, nhiệt độ và hiệu suất trong mọi bài toán thuộc dạng này.}
{\True H.P.t=λm khi mẫu ở nhiệt độ nóng chảy và bỏ qua phần khác.}
\loigiai{
\begin{itemchoice}
\item (Sai) Thời gian nóng chảy không phụ thuộc khối lượng. (Vì): Thời gian nóng chảy không phụ thuộc khối lượng. b) Đ: Khi giải dạng “Bài toán công suất, hiệu suất và thời gian nóng chảy”, cần xác định đúng trạng thái đầu, trạng thái cuối và quy ước dấu hoặc điều kiện áp dụng. c) S: Có thể bỏ qua mọi dữ kiện về khối lượng, nhiệt độ và hiệu suất trong mọi bài toán thuộc dạng này. d) Đ: H.P.t=λm khi mẫu ở nhiệt độ nóng chảy và bỏ qua phần khác. Kết luận dựa trên định nghĩa, điều kiện áp dụng và quy ước trong SGK KNTT
\item (Đúng) Khi giải dạng “Bài toán công suất, hiệu suất và thời gian nóng chảy”, cần xác định đúng trạng thái đầu, trạng thái cuối và quy ước dấu hoặc điều kiện áp dụng.
\item (Sai) Có thể bỏ qua mọi dữ kiện về khối lượng, nhiệt độ và hiệu suất trong mọi bài toán thuộc dạng này.
\item (Đúng) H.P.t=λm khi mẫu ở nhiệt độ nóng chảy và bỏ qua phần khác.
\end{itemchoice}
}
\end{ex}

% ===== Câu 6 =====
% ID: L12C1B6-05-TL
% Mức: NB
\begin{bt}
% ID: L12C1B6-05-TL
% Mức: NB
Trình bày cách nhận biết và phương pháp giải dạng “Bài toán công suất, hiệu suất và thời gian nóng chảy”. Phân tích nhận định sau và sửa lại nếu sai: “Thời gian nóng chảy không phụ thuộc khối lượng.”
\loigiai{
Trước hết xác định hiện tượng, trạng thái và các đại lượng đã cho. Nêu định nghĩa hoặc hệ thức phù hợp, đổi về đơn vị SI, áp dụng đúng điều kiện và kiểm tra ý nghĩa kết quả. Nhận định cần đối chiếu: H.P.t=λm khi mẫu ở nhiệt độ nóng chảy và bỏ qua phần khác. Vì vậy nhận định đã cho sai; phát biểu đúng là: H.P.t=λm khi mẫu ở nhiệt độ nóng chảy và bỏ qua phần khác.
}
\end{bt}

% ===== Câu 8 =====
% ID: L12C1B6-06-TN
% Mức: NB
\begin{ex}
% ID: L12C1B6-06-TN
% Mức: NB
Câu 5 thuộc dạng “Bài toán công suất, hiệu suất và thời gian nóng chảy”. Phát biểu nào sau đây đúng?
\choice
{Thời gian nóng chảy không phụ thuộc khối lượng.}
{Nhiệt lượng và nhiệt độ là hai đại lượng hoàn toàn giống nhau.}
{Mọi quá trình nhiệt học đều không phụ thuộc điều kiện thực hiện.}
{\True H.P.t=λm khi mẫu ở nhiệt độ nóng chảy và bỏ qua phần khác.}
\loigiai{
Phát biểu đúng là: H.P.t=λm khi mẫu ở nhiệt độ nóng chảy và bỏ qua phần khác. Các phương án còn lại trái với mô hình, định nghĩa hoặc hệ thức nhiệt học tương ứng.
}
\end{ex}

% ===== Câu 9 =====
% ID: L12C1B6-07-DS
% Mức: TH
\begin{ex}
% ID: L12C1B6-07-DS
% Mức: TH
Xét các phát biểu sau về dạng “Bài toán công suất, hiệu suất và thời gian nóng chảy” (bộ số 5).
\choiceTF
{Thời gian nóng chảy không phụ thuộc khối lượng.}
{\True Khi giải dạng “Bài toán công suất, hiệu suất và thời gian nóng chảy”, cần xác định đúng trạng thái đầu, trạng thái cuối và quy ước dấu hoặc điều kiện áp dụng.}
{Có thể bỏ qua mọi dữ kiện về khối lượng, nhiệt độ và hiệu suất trong mọi bài toán thuộc dạng này.}
{\True H.P.t=λm khi mẫu ở nhiệt độ nóng chảy và bỏ qua phần khác.}
\loigiai{
\begin{itemchoice}
\item (Sai) Thời gian nóng chảy không phụ thuộc khối lượng. (Vì): Thời gian nóng chảy không phụ thuộc khối lượng. b) Đ: Khi giải dạng “Bài toán công suất, hiệu suất và thời gian nóng chảy”, cần xác định đúng trạng thái đầu, trạng thái cuối và quy ước dấu hoặc điều kiện áp dụng. c) S: Có thể bỏ qua mọi dữ kiện về khối lượng, nhiệt độ và hiệu suất trong mọi bài toán thuộc dạng này. d) Đ: H.P.t=λm khi mẫu ở nhiệt độ nóng chảy và bỏ qua phần khác. Kết luận dựa trên định nghĩa, điều kiện áp dụng và quy ước trong SGK KNTT
\item (Đúng) Khi giải dạng “Bài toán công suất, hiệu suất và thời gian nóng chảy”, cần xác định đúng trạng thái đầu, trạng thái cuối và quy ước dấu hoặc điều kiện áp dụng.
\item (Sai) Có thể bỏ qua mọi dữ kiện về khối lượng, nhiệt độ và hiệu suất trong mọi bài toán thuộc dạng này.
\item (Đúng) H.P.t=λm khi mẫu ở nhiệt độ nóng chảy và bỏ qua phần khác.
\end{itemchoice}
}
\end{ex}

% ===== Câu 10 =====
% ID: L12C1B6-08-TL
% Mức: TH
\begin{bt}
% ID: L12C1B6-08-TL
% Mức: TH
Trình bày cách nhận biết và phương pháp giải dạng “Bài toán công suất, hiệu suất và thời gian nóng chảy”. Phân tích nhận định sau và sửa lại nếu sai: “H.P.t=λm khi mẫu ở nhiệt độ nóng chảy và bỏ qua phần khác.”
\loigiai{
Trước hết xác định hiện tượng, trạng thái và các đại lượng đã cho. Nêu định nghĩa hoặc hệ thức phù hợp, đổi về đơn vị SI, áp dụng đúng điều kiện và kiểm tra ý nghĩa kết quả. Nhận định cần đối chiếu: H.P.t=λm khi mẫu ở nhiệt độ nóng chảy và bỏ qua phần khác. Vì vậy nhận định đã cho đúng.
}
\end{bt}

% ===== Câu 12 =====
% ID: L12C1B6-09-TN
% Mức: NB
\begin{ex}
% ID: L12C1B6-09-TN
% Mức: NB
Câu 9 thuộc dạng “Bài toán công suất, hiệu suất và thời gian nóng chảy”. Phát biểu nào sau đây đúng?
\choice
{Thời gian nóng chảy không phụ thuộc khối lượng.}
{Nhiệt lượng và nhiệt độ là hai đại lượng hoàn toàn giống nhau.}
{Mọi quá trình nhiệt học đều không phụ thuộc điều kiện thực hiện.}
{\True H.P.t=λm khi mẫu ở nhiệt độ nóng chảy và bỏ qua phần khác.}
\loigiai{
Phát biểu đúng là: H.P.t=λm khi mẫu ở nhiệt độ nóng chảy và bỏ qua phần khác. Các phương án còn lại trái với mô hình, định nghĩa hoặc hệ thức nhiệt học tương ứng.
}
\end{ex}

% ===== Câu 13 =====
% ID: L12C1B6-10-DS
% Mức: VD
\begin{ex}
% ID: L12C1B6-10-DS
% Mức: VD
Xét các phát biểu sau về dạng “Bài toán công suất, hiệu suất và thời gian nóng chảy” (bộ số 2).
\choiceTF
{\True Khi giải dạng “Bài toán công suất, hiệu suất và thời gian nóng chảy”, cần xác định đúng trạng thái đầu, trạng thái cuối và quy ước dấu hoặc điều kiện áp dụng.}
{Có thể bỏ qua mọi dữ kiện về khối lượng, nhiệt độ và hiệu suất trong mọi bài toán thuộc dạng này.}
{\True H.P.t= $\lambda$ m khi mẫu ở nhiệt độ nóng chảy và bỏ qua phần khác.}
{Thời gian nóng chảy không phụ thuộc khối lượng.}
\loigiai{
\begin{itemchoice}
\item (Đúng) Khi giải dạng “Bài toán công suất, hiệu suất và thời gian nóng chảy”, cần xác định đúng trạng thái đầu, trạng thái cuối và quy ước dấu hoặc điều kiện áp dụng. (Vì): Khi giải các bài toán về công suất, hiệu suất và thời gian nóng chảy, việc xác định chính xác trạng thái ban đầu, trạng thái cuối và áp dụng đúng quy ước dấu hoặc điều kiện cho từng trường hợp là rất quan trọng để có kết quả chính xác
\item (Sai) Có thể bỏ qua mọi dữ kiện về khối lượng, nhiệt độ và hiệu suất trong mọi bài toán thuộc dạng này. (Vì): Khối lượng, nhiệt độ và hiệu suất là những yếu tố cốt lõi trong các bài toán này. Bỏ qua chúng sẽ dẫn đến sai lầm nghiêm trọng trong tính toán
\item (Đúng) H.P.t= $\lambda$ m khi mẫu ở nhiệt độ nóng chảy và bỏ qua phần khác. (Vì): Công thức $Q = \lambda m$ biểu diễn nhiệt lượng cần thiết để làm nóng chảy hoàn toàn một khối lượng $m$ của chất rắn ở nhiệt độ nóng chảy. Nếu $P$ là công suất tỏa nhiệt và $t$ là thời gian, thì $P \cdot t = Q = \lambda m$, với $\lambda$ là nhiệt nóng chảy riêng
\item (Sai) Thời gian nóng chảy không phụ thuộc khối lượng. (Vì): Thời gian nóng chảy phụ thuộc vào nhiệt lượng cần thiết để nóng chảy và công suất cung cấp. Nhiệt lượng cần thiết lại tỉ lệ thuận với khối lượng ($Q = \lambda m$), do đó thời gian nóng chảy sẽ phụ thuộc vào khối lượng
\end{itemchoice}
}
\end{ex}

% ===== Câu 14 =====
% ID: L12C1B6-19-TL
% Mức: VD
\begin{bt}
% ID: L12C1B6-19-TL
% Mức: VD
Trình bày cách nhận biết và phương pháp giải dạng “Bài toán công suất, hiệu suất và thời gian nóng chảy”. Phân tích nhận định sau và sửa lại nếu sai: “Thời gian nóng chảy không phụ thuộc khối lượng.”
\loigiai{
Trước hết xác định hiện tượng, trạng thái và các đại lượng đã cho. Nêu định nghĩa hoặc hệ thức phù hợp, đổi về đơn vị SI, áp dụng đúng điều kiện và kiểm tra ý nghĩa kết quả. Nhận định cần đối chiếu: H.P.t=λm khi mẫu ở nhiệt độ nóng chảy và bỏ qua phần khác. Vì vậy nhận định đã cho sai; phát biểu đúng là: H.P.t=λm khi mẫu ở nhiệt độ nóng chảy và bỏ qua phần khác.
}
\end{bt}

% ===== Câu 16 =====
% ID: L12C1B6-20-TN
% Mức: TH
\begin{ex}
% ID: L12C1B6-20-TN
% Mức: TH
Câu 2 thuộc dạng “Bài toán công suất, hiệu suất và thời gian nóng chảy”. Phát biểu nào sau đây đúng?
\choice
{Nhiệt lượng và nhiệt độ là hai đại lượng hoàn toàn giống nhau.}
{Mọi quá trình nhiệt học đều không phụ thuộc điều kiện thực hiện.}
{\True H.P.t=λm khi mẫu ở nhiệt độ nóng chảy và bỏ qua phần khác.}
{Thời gian nóng chảy không phụ thuộc khối lượng.}
\loigiai{
Phát biểu đúng là: H.P.t=λm khi mẫu ở nhiệt độ nóng chảy và bỏ qua phần khác. Các phương án còn lại trái với mô hình, định nghĩa hoặc hệ thức nhiệt học tương ứng.
}
\end{ex}

% ===== Câu 17 =====
% ID: L12C1B6-21-DS
% Mức: VDC
\begin{ex}
% ID: L12C1B6-21-DS
% Mức: VDC
Xét các phát biểu sau về dạng “Bài toán công suất, hiệu suất và thời gian nóng chảy” (bộ số 3).
\choiceTF
{Có thể bỏ qua mọi dữ kiện về khối lượng, nhiệt độ và hiệu suất trong mọi bài toán thuộc dạng này.}
{\True H.P.t=λm khi mẫu ở nhiệt độ nóng chảy và bỏ qua phần khác.}
{Thời gian nóng chảy không phụ thuộc khối lượng.}
{\True Khi giải dạng “Bài toán công suất, hiệu suất và thời gian nóng chảy”, cần xác định đúng trạng thái đầu, trạng thái cuối và quy ước dấu hoặc điều kiện áp dụng.}
\loigiai{
\begin{itemchoice}
\item (Sai) Có thể bỏ qua mọi dữ kiện về khối lượng, nhiệt độ và hiệu suất trong mọi bài toán thuộc dạng này. (Vì): Có thể bỏ qua mọi dữ kiện về khối lượng, nhiệt độ và hiệu suất trong mọi bài toán thuộc dạng này. b) Đ: H.P.t=λm khi mẫu ở nhiệt độ nóng chảy và bỏ qua phần khác. c) S: Thời gian nóng chảy không phụ thuộc khối lượng. d) Đ: Khi giải dạng “Bài toán công suất, hiệu suất và thời gian nóng chảy”, cần xác định đúng trạng thái đầu, trạng thái cuối và quy ước dấu hoặc điều kiện áp dụng. Kết luận dựa trên định nghĩa, điều kiện áp dụng và quy ước trong SGK KNTT
\item (Đúng) H.P.t=λm khi mẫu ở nhiệt độ nóng chảy và bỏ qua phần khác.
\item (Sai) Thời gian nóng chảy không phụ thuộc khối lượng.
\item (Đúng) Khi giải dạng “Bài toán công suất, hiệu suất và thời gian nóng chảy”, cần xác định đúng trạng thái đầu, trạng thái cuối và quy ước dấu hoặc điều kiện áp dụng.
\end{itemchoice}
}
\end{ex}

% ===== Câu 18 =====
% ID: L12C1B6-37-TL
% Mức: VDC
\begin{bt}
% ID: L12C1B6-37-TL
% Mức: VDC
Trình bày cách nhận biết và phương pháp giải dạng “Bài toán công suất, hiệu suất và thời gian nóng chảy”. Phân tích nhận định sau và sửa lại nếu sai: “H.P.t=λm khi mẫu ở nhiệt độ nóng chảy và bỏ qua phần khác.”
\loigiai{
Trước hết xác định hiện tượng, trạng thái và các đại lượng đã cho. Nêu định nghĩa hoặc hệ thức phù hợp, đổi về đơn vị SI, áp dụng đúng điều kiện và kiểm tra ý nghĩa kết quả. Nhận định cần đối chiếu: H.P.t=λm khi mẫu ở nhiệt độ nóng chảy và bỏ qua phần khác. Vì vậy nhận định đã cho đúng.
}
\end{bt}

% ===== Câu 20 =====
% ID: L12C1B6-38-TN
% Mức: TH
\begin{ex}
% ID: L12C1B6-38-TN
% Mức: TH
Câu 6 thuộc dạng “Bài toán công suất, hiệu suất và thời gian nóng chảy”. Phát biểu nào sau đây đúng?
\choice
{Nhiệt lượng và nhiệt độ là hai đại lượng hoàn toàn giống nhau.}
{Mọi quá trình nhiệt học đều không phụ thuộc điều kiện thực hiện.}
{\True H.P.t=λm khi mẫu ở nhiệt độ nóng chảy và bỏ qua phần khác.}
{Thời gian nóng chảy không phụ thuộc khối lượng.}
\loigiai{
Phát biểu đúng là: H.P.t=λm khi mẫu ở nhiệt độ nóng chảy và bỏ qua phần khác. Các phương án còn lại trái với mô hình, định nghĩa hoặc hệ thức nhiệt học tương ứng.
}
\end{ex}

% ===== Câu 22 =====
% ID: L12C1B6-39-TL
% Mức: VDC
\begin{bt}
% ID: L12C1B6-39-TL
% Mức: VDC
Trình bày cách nhận biết và phương pháp giải dạng “Bài toán công suất, hiệu suất và thời gian nóng chảy”. Phân tích nhận định sau và sửa lại nếu sai: “H.P.t=λm khi mẫu ở nhiệt độ nóng chảy và bỏ qua phần khác.”
\loigiai{
Trước hết xác định hiện tượng, trạng thái và các đại lượng đã cho. Nêu định nghĩa hoặc hệ thức phù hợp, đổi về đơn vị SI, áp dụng đúng điều kiện và kiểm tra ý nghĩa kết quả. Nhận định cần đối chiếu: H.P.t=λm khi mẫu ở nhiệt độ nóng chảy và bỏ qua phần khác. Vì vậy nhận định đã cho đúng.
}
\end{bt}

% ===== Câu 24 =====
% ID: L12C1B6-40-TN
% Mức: TH
\begin{ex}
% ID: L12C1B6-40-TN
% Mức: TH
Câu 10 thuộc dạng “Bài toán công suất, hiệu suất và thời gian nóng chảy”. Phát biểu nào sau đây đúng?
\choice
{Nhiệt lượng và nhiệt độ là hai đại lượng hoàn toàn giống nhau.}
{Mọi quá trình nhiệt học đều không phụ thuộc điều kiện thực hiện.}
{\True H.P.t=λm khi mẫu ở nhiệt độ nóng chảy và bỏ qua phần khác.}
{Thời gian nóng chảy không phụ thuộc khối lượng.}
\loigiai{
Phát biểu đúng là: H.P.t=λm khi mẫu ở nhiệt độ nóng chảy và bỏ qua phần khác. Các phương án còn lại trái với mô hình, định nghĩa hoặc hệ thức nhiệt học tương ứng.
}
\end{ex}

% ===== Câu 28 =====
% ID: L12C1B6-41-TN
% Mức: VD
\begin{ex}
% ID: L12C1B6-41-TN
% Mức: VD
Câu 3 thuộc dạng “Bài toán công suất, hiệu suất và thời gian nóng chảy”. Phát biểu nào sau đây đúng?
\choice
{Mọi quá trình nhiệt học đều không phụ thuộc điều kiện thực hiện.}
{\True H.P.t=λm khi mẫu ở nhiệt độ nóng chảy và bỏ qua phần khác.}
{Thời gian nóng chảy không phụ thuộc khối lượng.}
{Nhiệt lượng và nhiệt độ là hai đại lượng hoàn toàn giống nhau.}
\loigiai{
Phát biểu đúng là: H.P.t=λm khi mẫu ở nhiệt độ nóng chảy và bỏ qua phần khác. Các phương án còn lại trái với mô hình, định nghĩa hoặc hệ thức nhiệt học tương ứng.
}
\end{ex}

% ===== Câu 32 =====
% ID: L12C1B6-42-TN
% Mức: VD
\begin{ex}
% ID: L12C1B6-42-TN
% Mức: VD
Câu 7 thuộc dạng “Bài toán công suất, hiệu suất và thời gian nóng chảy”. Phát biểu nào sau đây đúng?
\choice
{Mọi quá trình nhiệt học đều không phụ thuộc điều kiện thực hiện.}
{\True H.P.t=λm khi mẫu ở nhiệt độ nóng chảy và bỏ qua phần khác.}
{Thời gian nóng chảy không phụ thuộc khối lượng.}
{Nhiệt lượng và nhiệt độ là hai đại lượng hoàn toàn giống nhau.}
\loigiai{
Phát biểu đúng là: H.P.t=λm khi mẫu ở nhiệt độ nóng chảy và bỏ qua phần khác. Các phương án còn lại trái với mô hình, định nghĩa hoặc hệ thức nhiệt học tương ứng.
}
\end{ex}

% ===== Câu 36 =====
% ID: L12C1B6-43-TN
% Mức: VDC
\begin{ex}
% ID: L12C1B6-43-TN
% Mức: VDC
Câu 4 thuộc dạng “Bài toán công suất, hiệu suất và thời gian nóng chảy”. Phát biểu nào sau đây đúng?
\choice
{\True H.P.t=λm khi mẫu ở nhiệt độ nóng chảy và bỏ qua phần khác.}
{Thời gian nóng chảy không phụ thuộc khối lượng.}
{Nhiệt lượng và nhiệt độ là hai đại lượng hoàn toàn giống nhau.}
{Mọi quá trình nhiệt học đều không phụ thuộc điều kiện thực hiện.}
\loigiai{
Phát biểu đúng là: H.P.t=λm khi mẫu ở nhiệt độ nóng chảy và bỏ qua phần khác. Các phương án còn lại trái với mô hình, định nghĩa hoặc hệ thức nhiệt học tương ứng.
}
\end{ex}

% ===== Câu 40 =====
% ID: L12C1B6-44-TN
% Mức: VDC
\begin{ex}
% ID: L12C1B6-44-TN
% Mức: VDC
Câu 8 thuộc dạng “Bài toán công suất, hiệu suất và thời gian nóng chảy”. Phát biểu nào sau đây đúng?
\choice
{\True H.P.t=λm khi mẫu ở nhiệt độ nóng chảy và bỏ qua phần khác.}
{Thời gian nóng chảy không phụ thuộc khối lượng.}
{Nhiệt lượng và nhiệt độ là hai đại lượng hoàn toàn giống nhau.}
{Mọi quá trình nhiệt học đều không phụ thuộc điều kiện thực hiện.}
\loigiai{
Phát biểu đúng là: H.P.t=λm khi mẫu ở nhiệt độ nóng chảy và bỏ qua phần khác. Các phương án còn lại trái với mô hình, định nghĩa hoặc hệ thức nhiệt học tương ứng.
}
\end{ex}

% ===== Câu 56 =====
% ID: L12C1B6-45-TLN
% Mức: NB
\begin{ex}
% ID: L12C1B6-45-TLN
% Mức: NB
Cho bốn nhận định: (1) H.P.t=λm khi mẫu ở nhiệt độ nóng chảy và bỏ qua phần khác. (2) Thời gian nóng chảy không phụ thuộc khối lượng. (3) Cần kiểm tra đơn vị và điều kiện áp dụng khi xử lí dạng “Bài toán công suất, hiệu suất và thời gian nóng chảy”. (4) Có thể thay tùy ý nhiệt độ Celsius cho nhiệt độ tuyệt đối trong mọi công thức. Có bao nhiêu nhận định đúng? Trả lời bằng một số nguyên.
\shortans{2}
\loigiai{
Nhận định (1) và (3) đúng; (2) và (4) sai. Vậy có 2 nhận định đúng.
}
\end{ex}

% ===== Câu 60 =====
% ID: L12C1B6-46-TLN
% Mức: TH
\begin{ex}
% ID: L12C1B6-46-TLN
% Mức: TH
Cho bốn nhận định: (1) H.P.t=λm khi mẫu ở nhiệt độ nóng chảy và bỏ qua phần khác. (2) Thời gian nóng chảy không phụ thuộc khối lượng. (3) Cần kiểm tra đơn vị và điều kiện áp dụng khi xử lí dạng “Bài toán công suất, hiệu suất và thời gian nóng chảy”. (4) Có thể thay tùy ý nhiệt độ Celsius cho nhiệt độ tuyệt đối trong mọi công thức. Có bao nhiêu nhận định đúng? Trả lời bằng một số nguyên.
\shortans{2}
\loigiai{
Nhận định (1) và (3) đúng; (2) và (4) sai. Vậy có 2 nhận định đúng.
}
\end{ex}

% ===== Câu 64 =====
% ID: L12C1B6-47-TLN
% Mức: VD
\begin{ex}
% ID: L12C1B6-47-TLN
% Mức: VD
Cho bốn nhận định: (1) H.P.t=λm khi mẫu ở nhiệt độ nóng chảy và bỏ qua phần khác. (2) Thời gian nóng chảy không phụ thuộc khối lượng. (3) Cần kiểm tra đơn vị và điều kiện áp dụng khi xử lí dạng “Bài toán công suất, hiệu suất và thời gian nóng chảy”. (4) Có thể thay tùy ý nhiệt độ Celsius cho nhiệt độ tuyệt đối trong mọi công thức. Có bao nhiêu nhận định đúng? Trả lời bằng một số nguyên.
\shortans{2}
\loigiai{
Nhận định (1) và (3) đúng; (2) và (4) sai. Vậy có 2 nhận định đúng.
}
\end{ex}

% ===== Câu 68 =====
% ID: L12C1B6-48-TLN
% Mức: VD
\begin{ex}
% ID: L12C1B6-48-TLN
% Mức: VD
Cho bốn nhận định: (1) H.P.t=λm khi mẫu ở nhiệt độ nóng chảy và bỏ qua phần khác. (2) Thời gian nóng chảy không phụ thuộc khối lượng. (3) Cần kiểm tra đơn vị và điều kiện áp dụng khi xử lí dạng “Bài toán công suất, hiệu suất và thời gian nóng chảy”. (4) Có thể thay tùy ý nhiệt độ Celsius cho nhiệt độ tuyệt đối trong mọi công thức. Có bao nhiêu nhận định đúng? Trả lời bằng một số nguyên.
\shortans{2}
\loigiai{
Nhận định (1) và (3) đúng; (2) và (4) sai. Vậy có 2 nhận định đúng.
}
\end{ex}

% ===== Câu 72 =====
% ID: L12C1B6-49-TLN
% Mức: VDC
\begin{ex}
% ID: L12C1B6-49-TLN
% Mức: VDC
Cho bốn nhận định: (1) H.P.t=λm khi mẫu ở nhiệt độ nóng chảy và bỏ qua phần khác. (2) Thời gian nóng chảy không phụ thuộc khối lượng. (3) Cần kiểm tra đơn vị và điều kiện áp dụng khi xử lí dạng “Bài toán công suất, hiệu suất và thời gian nóng chảy”. (4) Có thể thay tùy ý nhiệt độ Celsius cho nhiệt độ tuyệt đối trong mọi công thức. Có bao nhiêu nhận định đúng? Trả lời bằng một số nguyên.
\shortans{2}
\loigiai{
Nhận định (1) và (3) đúng; (2) và (4) sai. Vậy có 2 nhận định đúng.
}
\end{ex}

% ===== Câu 76 =====
% ID: L12C1B6-50-TLN
% Mức: VDC
\begin{ex}
% ID: L12C1B6-50-TLN
% Mức: VDC
Cho bốn nhận định: (1) H.P.t=λm khi mẫu ở nhiệt độ nóng chảy và bỏ qua phần khác. (2) Thời gian nóng chảy không phụ thuộc khối lượng. (3) Cần kiểm tra đơn vị và điều kiện áp dụng khi xử lí dạng “Bài toán công suất, hiệu suất và thời gian nóng chảy”. (4) Có thể thay tùy ý nhiệt độ Celsius cho nhiệt độ tuyệt đối trong mọi công thức. Có bao nhiêu nhận định đúng? Trả lời bằng một số nguyên.
\shortans{2}
\loigiai{
Nhận định (1) và (3) đúng; (2) và (4) sai. Vậy có 2 nhận định đúng.
}
\end{ex}
